% Generated draft: requires mathematical and source review.
\documentclass[11pt]{article}
\usepackage[margin=25mm]{geometry}
\usepackage{fontspec}
\setmainfont{Latin Modern Roman}
\usepackage{amsmath,amssymb,mathtools,unicode-math}
\setmathfont{Latin Modern Math}
\usepackage{enumitem,needspace,xurl,hyperref}
\setlength{\parindent}{0pt}
\setlength{\parskip}{5pt}
\setlength{\emergencystretch}{3em}
\newcommand{\N}{\mathbb{N}}
\newcommand{\Z}{\mathbb{Z}}
\newcommand{\Q}{\mathbb{Q}}
\newcommand{\R}{\mathbb{R}}
\newcommand{\C}{\mathbb{C}}
\newcommand{\F}{\mathbb{F}}
\newcommand{\A}{\mathbb{A}}
\newcommand{\PP}{\mathbb{P}}
\newcommand{\E}{\mathbb{E}}
\newcommand{\eps}{\varepsilon}
\newcommand{\dd}{\,\mathrm{d}}
\begin{document}
\section*{UnsolvedMath ID: 20000318}
\textbf{Source problem code:} catalogue entry
\subsection{Definitions and mathematical statement}
Let $X_n\subset\PP^4$ be a degree-$n$ hypersurface with a finite set $\Sigma=\operatorname{Sing}X_n$ of ordinary double points and no other singularities. In the attached aid the field is $\C$; the extracted original does not state it. An ordinary double point is a hypersurface singularity with nondegenerate quadratic part in the intended characteristic-zero setting. Write $\operatorname{Cl}(X_n)$ for Weil divisors modulo linear equivalence. The source uses ``factorial'' to mean $\operatorname{Cl}(X_n)=\Z$, with the hyperplane class understood as generator.

\textbf{Original questions (Cheltsov).} Determine when $X_n$ is factorial. In particular, does
\[|\Sigma|<(n-1)^2\quad\Longrightarrow\quad\operatorname{Cl}(X_n)=\Z?\]
Also, for a generic linear projection $\phi:\PP^4\dashrightarrow\PP^2$, is it true that, for every integer $k\ge1$ and every plane curve $B$ of degree $k$,
\[|\phi(\Sigma)\cap B|\le k(n-1)?\]
The projection should be defined and injective on the finite set $\Sigma$, as occurs for a sufficiently general center. The extracted projection question states no cardinality hypothesis on $\Sigma$; this must not be silently replaced by the more restricted version discussed by the aid.

The original attributes to Clemens a criterion using independent conditions on degree-$n$ forms. The aid instead gives degree $2n-5$:
\[X_n\text{ factorial}\quad\Longleftrightarrow\quad H^0(\PP^4,\mathcal O(2n-5))\longrightarrow H^0(\Sigma,\mathcal O_\Sigma(2n-5))\text{ is surjective}.\]
This correction is explicitly attributed to the aid; the inconsistent extracted degree is retained as a source issue. One may measure the indicated defect by $\delta(X_n)=h^1(\PP^4,\mathcal I_\Sigma(2n-5))$.

The displayed source example, in coordinates $x,y,z,w,v$, is
\[X_n=\{xf_{n-1}+yg_{n-1}=0\},\]
where $f_{n-1},g_{n-1}$ are generic homogeneous forms. It reports $\operatorname{Cl}(X_n)=\Z\oplus\Z$. In the intended $n\ge2$ nodal example the hypersurface contains the plane $x=y=0$; its node set in the equality case is of complete-intersection type $(1,1,n-1,n-1)$.

Two extracted ``known facts'' require caution. The first has a damaged coefficient printed as $2$, an intervening comparison sign, and $3$ before $(n-1)^2$, so its numerical hypothesis cannot be certified from this record. The second says that every smooth codimension-one subvariety $S\subset X_n$ is Cartier. Its direction or omitted qualification needs checking: taken literally it conflicts with the role of the smooth contained plane in the displayed nonfactorial example. Neither damaged claim is accepted here as a verified theorem.

\textbf{Unreviewed research-aid claims.} The aid attributes to Cheltsov the sharp affirmative threshold over $\C$: every nodal degree-$n$ hypersurface with at most $(n-1)^2-1$ nodes is factorial. It says that under the exact Cheltsov--Park cardinality hypothesis $|\Sigma|\le(n-1)^2-1$, the full generic-projection incidence assertion was not found resolved, while the required bound was proved for $k\le\sqrt{n-1}$. This is a report about that restricted formulation and literature search, not an assertion that the unrestricted extracted question has the same status.

Finally, the aid claims that every nonfactorial complex nodal degree-$n$ hypersurface with exactly $(n-1)^2$ nodes has reduced node scheme a complete intersection of type $(1,1,n-1,n-1)$, defect $\delta(X_n)=1$, and $\operatorname{Cl}(X_n)\cong\Z^2$, not just generic hypersurfaces of the displayed equation. Its necessary degree range, reference and full proof are not supplied. These assertions require independent mathematical and source review.
\subsection{Short English statement}
Cheltsov asks for factoriality and a generic-projection incidence bound for nodal threefold hypersurfaces. The aid reports the sharp node-count threshold and a defect-one equality claim. The draft flags the conflicting criterion degrees, damaged historical bound, questionable smooth-divisor statement, and the cardinality hypothesis present only in the aid's projection discussion.
\subsection{Sources}
\begin{itemize}
\item Local UnsolvedMath registry, permanent ID 20000318, source code catalogue entry; \texttt{aim-algebraic-geometry-notes.json}, notes[317], tag \texttt{problem}.
\item AIM workshop \emph{Numerical invariants of singularities and higher-dimensional algebraic varieties}, Problem 8 (Cheltsov): \url{https://aimath.org/WWN/singularvariety/problems_final.pdf}. The original names Clemens. The supplied status is \texttt{reconstructed\_unverified}; its verification field refers to original TeX but supplies no complete repaired transcription. No PDF or TeX source was consulted here.
\item Attached research attempt 1 names Cheltsov and Cheltsov--Park without bibliographic locators; original literature field empty. Result \texttt{partial\_result}, status \texttt{partially\_solved}, rigor label \texttt{proved}, novelty confidence low. It is a machine-generated, non-peer-reviewed aid requiring independent verification before citation.
\end{itemize}
\end{document}
